\documentclass[11pt]{article}

% Local NIH-style Specific Aims scaffold.
% This is NOT an official NIH template.
% Verify the current FOA and NIH application guide before submission.

\usepackage[margin=0.5in]{geometry}
\usepackage{times}
\usepackage{hyperref}

\newcommand{\todo}[1]{\textbf{TODO: #1}}

\begin{document}

\section*{Specific Aims}

\todo{Verify current NIH FOA, page limits, font/margin requirements, and required components.}

\paragraph{Significance.}
\todo{State the biomedical/clinical problem only if applicable. Use citations and claim-auditor for health-related claims.}

\paragraph{Gap.}
\todo{State the specific gap. Do not invent clinical need or preliminary evidence.}

\paragraph{Objective.}
\todo{State the objective of this proposal.}

\paragraph{Central Hypothesis.}
\todo{State a testable hypothesis. Mark as hypothesis, not finding.}

\paragraph{Aim 1.}
\todo{Describe aim, approach, expected outcome, and risk mitigation.}

\paragraph{Aim 2.}
\todo{Describe aim, approach, expected outcome, and risk mitigation.}

\paragraph{Aim 3.}
\todo{Optional. Include only if justified by scope.}

\paragraph{Impact.}
\todo{State realistic impact proportional to evidence and proposal scope.}

\end{document}
